% ===================================================================
%  Basis für Blanko-Notenpapier (A5) – wird von den vier Vorlagen
%  vorlage_notenpapier_a5_*.tex per \input geladen, nicht allein kompilieren.
%
%  Die Vorlage setzt vorher zwei Schalter:
%    \querfalse / \quertrue               Hoch- oder Querformat
%    \schluesseltrue / \schluesselfalse   mit oder ohne Violinschlüssel
%  Reines TikZ; der Violinschlüssel ist Zeichen 71 der MusiXTeX-Schrift
%  musix20 (nur als Schrift geladen, kein \usepackage{musixtex}).
% ===================================================================
\ifquer
  \usepackage[a5paper,landscape,margin=12mm]{geometry}
\else
  \usepackage[a5paper,margin=12mm]{geometry}
\fi
\usepackage{tikz}
\pagestyle{empty}
\setlength{\parindent}{0pt}

% ---- Einstellungen ----------------------------------------------------
\newlength{\Linienabstand}  \setlength{\Linienabstand}{2.5mm}  % Abstand zweier Notenlinien
\newlength{\Systemabstand}  \setlength{\Systemabstand}{11mm}   % Lücke zwischen zwei Systemen
\newcommand{\Linienstaerke}{0.4pt}

% Schlüsselgröße passt sich dem Linienabstand an: musix20 gehört zu 5pt
% Linienabstand bei 20pt Schriftgröße, also Schriftgröße = 4 x Linienabstand.
\font\notenschrift=musix20 at \dimexpr 4\Linienabstand\relax

\begin{document}
\begin{tikzpicture}[x=1mm,y=1mm]
  \pgfmathsetmacro{\L}{\Linienabstand/1mm}
  \pgfmathsetmacro{\G}{\Systemabstand/1mm}
  \pgfmathsetmacro{\B}{\linewidth/1mm}
  \pgfmathsetmacro{\H}{\textheight/1mm}
  % so viele Systeme, wie auf die Seite passen
  \pgfmathtruncatemacro{\n}{floor((\H+\G)/(4*\L+\G))}
  % Restplatz gleichmäßig oben und unten verteilen
  \pgfmathsetmacro{\oben}{(\H-(\n*4*\L+(\n-1)*\G))/2}
  % feste Zeichenfläche = Satzspiegel (Schlüssel dürfen in den Rand ragen)
  \path[use as bounding box] (0,0) rectangle (\B,-\H);
  \foreach \s in {1,...,\n} {
    \pgfmathsetmacro{\y}{-\oben-(\s-1)*(4*\L+\G)}  % oberste Linie des Systems
    \foreach \l in {0,...,4}
      \draw[line width=\Linienstaerke] (0,\y-\l*\L) -- (\B,\y-\l*\L);
    \ifschluessel
      % Grundlinie des Schlüssels auf der g'-Linie (2. Linie von unten)
      \node[anchor=base west, inner sep=0] at (0.8,\y-3*\L) {\notenschrift\char71};
    \fi
  }
\end{tikzpicture}
\end{document}
